\subsection{Involutive algebras}

DEFINITION 1. --- Let A be an algebra over C. An involution on A is a map \(x \mapsto x^{*}\) from A to A such that:

\[
\begin{array}{c} (x ^ {*}) ^ {*} = x, \quad (x + y) ^ {*} = x ^ {*} + y ^ {*}, \quad (\lambda x) ^ {*} = \overline {{\lambda}} x ^ {*} \\ (x y) ^ {*} = y ^ {*} x ^ {*} \end{array}
\]

for all \(x, y \in A\) and \(\lambda \in C\) . An algebra over C equipped with an involution is called an involutive algebra.

An involution on A is in particular an isomorphism of the ring A onto the opposite ring A°.

Let A be an involutive algebra. One says that \(x^{*}\) is the adjoint of x. A subset of A stable under the involution is said to be self-adjoint. If A has a unit element e, one has \(e^{*} = e\) ; one says that (A, e) is a unital involutive algebra. An element u of a unital involutive algebra is said to be unitary if \(uu^{*} = u^{*}u = e\) , in other words if u is invertible and its inverse is \(u^{*}\) .

An element \(x \in A\) is said to be hermitian if \(x = x^{*}\) and normal if \(xx^{*} = x^{*}x\) . This terminology generalizes that of A, IX, § 7, no. 3. Every hermitian element is normal, and every unitary element is normal. The set \(A_{h}\) of hermitian elements of A is a real vector subspace of A. If x and y are hermitian and commute, one has \((xy)^{*} = y^{*}x^{*} = yx = xy\) , hence xy is hermitian. For every \(x \in A\) , the elements \(xx^{*}\) and \(x^{*}x\) of A are hermitian.

If A = C equipped with the involution \(z \mapsto \overline{z}\) , one has \(A_{h} = R\) .

Lemma 2. --- Let A be an involutive algebra and \(x \in \mathrm{A}\) . The elements

\[
x _ {1} = \frac {1}{2} (x + x ^ {*}), \qquad x _ {2} = \frac {1}{2 i} (x - x ^ {*})
\]

are hermitian and satisfy \(x = x_{1} + ix_{2}\) . If \(x = y_{1} + iy_{2}\), where \(y_{1}\) and \(y_{2}\) are hermitian, then \(x_{1} = y_{1}\) and \(x_{2} = y_{2}\) . Moreover, the element x is normal if and only if \(x_{1}\) and \(x_{2}\) commute.

The first two assertions follow from lemma 1 of I, p.~95. One computes that \(xx^{*} - x^{*}x = 2i(x_{2}x_{1} - x_{1}x_{2})\) , hence x is normal if and only if \(x_{1}\) and \(x_{2}\) commute.

Let A be a unital involutive algebra. For \(x \in A\) to be invertible, it is necessary and sufficient that \(x^{*}\) be so, and one then has \((x^{*})^{-1} = (x^{-1})^{*}\) . Since \((x - \lambda e)^{*} = x^{*} - \overline{\lambda}e\) for all \(\lambda \in C\) , one deduces that \(\mathrm{Sp}_{\mathrm{A}}(x^{*}) = \overline{\mathrm{Sp}_{\mathrm{A}}(x)}\) .

Let A be an involutive algebra and \(\widetilde{\mathrm{A}}\) the algebra obtained from A by adjoining a unit element. There exists on \(\widetilde{\mathrm{A}}\) a unique involution extending that of A, given by \((\lambda, x)^{*} = (\overline{\lambda}, x^{*})\) for \(\lambda \in \mathbf{C}\) and \(x \in \mathrm{A}\) . If \(x \in \mathrm{A}\) , one has \(\mathrm{Sp}_{\mathrm{A}}'(x^{*}) = \overline{\mathrm{Sp}_{\mathrm{A}}'(x)}\) .

Let A and B be involutive algebras. A morphism from A to B is an algebra morphism \(\varphi\) from A to B such that \(\varphi(x^{*}) = \varphi(x)^{*}\) for all \(x\) in A. The identity map of A is a morphism of involutive algebras. If C is an involutive algebra and \(\pi: B \to C\) a morphism of involutive algebras, then \(\pi \circ \varphi\) is a morphism of involutive algebras. If \(\varphi\) is an isomorphism of algebras, then \(\varphi^{-1}\) is a morphism of involutive algebras, and one says that \(\varphi\) is an isomorphism of involutive algebras.

By prop. 1 of I, p.~95, if A and B are involutive algebras, an algebra morphism \(\varphi\) from A to B is a morphism of involutive algebras if and only if \(\varphi(A_h) \subset B_h\) . One calls an involutive subalgebra of A a self-adjoint subalgebra. The center of A is an involutive subalgebra. If \(A_1\) is a self-adjoint two-sided ideal of A, the involution of A defines by passage to the quotient an involution on the algebra \(A/A_1\) , and the canonical map from A onto \(A/A_1\) is a morphism of involutive algebras.

Let A be an involutive algebra. The radical of A is equal to the radical of the opposite algebra (A, VIII, p.~431, prop. 7), and is therefore self-adjoint.

Let A be an involutive algebra. If M ⊂ A is self-adjoint, its commutant M' is an involutive subalgebra of A. If \(x \in A\) , the bicommutant of \(\{x, x^{*}\}\) is an involutive subalgebra containing x and \(x^{*}\) and this subalgebra is commutative if and only if x is normal (no. 5 of I, p. 5).

Remark. — Let A be an involutive algebra and let B be a maximal commutative involutive subalgebra of A. Then B is a maximal commutative subalgebra. In particular, if A is unital, then the algebra B is full.

Indeed, let \(x \in A\) be an element commuting with B. Then \(x^{*}\) commutes with B. Write \(x = x_{1} + ix_{2}\) with \(x_{1}\) and \(x_{2}\) hermitian; the elements \(x_{1}\) and \(x_{2}\) commute with B (lemma 2). The subalgebra of A generated by B and \(x_{1}\) is therefore commutative and involutive. Consequently, it is equal to B, so that \(x_{1} \in B\) . Similarly, we have \(x_{2} \in B\) , and hence finally \(x \in B\) .

Let A be an involutive algebra. If \(f\) is a linear form on A, the map \(x \mapsto \overline{f(x^*)}\) on A is a linear form on A, which is denoted by \(f^*\) . The map \(f \mapsto f^*\) is a semilinear involution on A'. We say that \(f\) is Hermitian if \(f = f^*\). By lemma 1 of I, p.~95, every linear form \(f\) on A has a unique representation \(f = f_1 + if_2\) where \(f_1\) and \(f_2\) are Hermitian, namely \(f_1 = \frac{1}{2}(f + f^*)\) and \(f_2 = \frac{1}{2i}(f - f^*)\) .

For a linear form \(f\) to be Hermitian, it is necessary and sufficient that the restriction of \(f\) to \(A_h\) be real-valued (proposition 1 of I, p.~95). The map \(f \mapsto f|A_h\) is an isomorphism of the real vector space of Hermitian forms on the dual vector space of the real vector space \(A_h\) .

In particular, we denote by \(\mathsf{X}^{\prime}(\mathsf{A})_{h}\) (resp. \(\mathsf{X}(\mathsf{A})_{h}\) ) the set of Hermitian characters of A (resp. the set of nonzero Hermitian characters of A). A character is therefore Hermitian if its restriction to \(A_{h}\) takes real values.

If A is commutative and if \(\chi\) is a character of A, then \(\chi^{*}\) is a character of A, and the map \(\chi \mapsto \chi^{*}\) is a homeomorphism from \(\mathsf{X}'(\mathsf{A})\) onto \(\mathsf{X}'(\mathsf{A})\) .

Lemma 3. --- Let A be a commutative involutive algebra. The Gelfand transform of A to \(\mathcal{C}_{0}(\mathsf{X}(\mathsf{A}))\) is a morphism of involutive algebras if and only if every character of A is Hermitian.

Indeed, saying that \(G_{A}\) is a morphism of involutive algebras amounts to saying that, for all \(x \in A\) and \(\chi \in \mathsf{X}(A)\) , one has

\[
\chi (x ^ {*}) = \mathcal {G} _ {\mathrm{A}} (x ^ {*}) (\chi) = \overline {{\mathcal {G} _ {\mathrm{A}} (x) (\chi)}} = \overline {{\chi (x)}},
\]

that is, every \(\chi\) is Hermitian.

Examples. --- 1) Let A be the algebra of complex-valued functions on a set X. The map \(f \mapsto \overline{f}\) is an involution in A. The subalgebra of bounded functions in X is an involutive subalgebra of A. If X is a locally compact topological space, the subalgebras \(\mathcal{C}(X), \mathcal{C}_b(X), \mathcal{C}_0(X)\) and \(\mathcal{K}(X)\) are involutive subalgebras of A.

\begin{enumerate}
\def\labelenumi{\arabic{enumi})}
\setcounter{enumi}{1}
\item
  Let X be a locally compact topological space and \(\mu\) a positive measure on X. The map \(f \mapsto \overline{f}\) is an involution on the algebra \(\mathcal{L}^{\infty}(X, \mu)\) ; it induces, by passing to the quotient, an involution on the unital algebra \(\mathrm{L}^{\infty}(X, \mu)\) .
\item
  Let E be a complex Hilbert space. On the Banach algebra \(\mathcal{L}(\mathrm{E})\) , the map \(x \mapsto x^{*}\) (EVT, V, p.~37, prop. 1) is an involution.
\item
  Let G be a locally compact group. Let \(\mathcal{M}^{1}(G)\) be the Banach algebra of bounded complex measures on G (example 6 of I, p.~19).
\end{enumerate}

The map \(x \mapsto x^{-1}\) from G to G maps every measure \(\mu \in \mathcal{M}^{1}(G)\) into a measure \(\check{\mu} \in \mathcal{M}^{1}(G)\) (INT, VII, p.~12, formula (13)). We denote by \(\mu^{*}\) the complex conjugate measure of \(\check{\mu}\) . The map \(\mu \mapsto \check{\mu}\) is an isometric isomorphism from the Banach algebra \(\mathcal{M}^{1}(G)\) onto the Banach algebra \(\mathcal{M}^{1}(G^{\circ})\) (INT, VIII, §3, no. 1, cor. of prop. 7); hence \(\mu \mapsto \mu^{*}\) is an isometric involution of the Banach algebra \(\mathcal{M}^{1}(G)\) .

The set A of bounded measures admitting a density with respect to a Haar measure is a closed subalgebra of \(\mathcal{M}^{1}(\mathrm{G})\) stable under the involution (cf. INT, VIII, §4, no. 5); it does not depend on the choice of a Haar measure.

Let \(\nu\) be a left Haar measure on G and denote by \(\Delta\) the module of G. We equip \(\mathrm{L}^1 (\mathrm{G},\nu)\) with the product \((f,g)\mapsto f*\nu g\) and with the involution \(f\mapsto f^{*} = \widetilde{f}\cdot \Delta^{-1}\) , where \(\widetilde{f} (x) = \overline{f(x^{-1})}\) for all \(x\in \mathbf{G}\). Then the map \(f\mapsto f\cdot \nu\) is an isomorphism of the involutive algebra \(\mathrm{L}^1 (\mathrm{G},\nu)\) onto A. This isomorphism is isometric. In particular, \(\mathrm{L}^1 (\mathrm{G},\nu)\) is identified with an involutive subalgebra of \(\mathcal{M}^1 (\mathrm{G})\) .

\begin{enumerate}
\def\labelenumi{\arabic{enumi})}
\setcounter{enumi}{4}
\tightlist
\item
  Let U be an open subset of C stable under complex conjugation. Consider the algebra \(\mathcal{O}(\mathrm{U})\) of complex-valued holomorphic functions on U. For every function \(f \in \mathcal{O}(\mathrm{U})\) , the map \(f^{*}: z \mapsto \overline{f(\overline{z})}\) is a holomorphic function on U. The map \(f \mapsto f^{*}\) is an involution on \(\mathcal{O}(\mathrm{U})\) .
\end{enumerate}

Similarly, let S be a compact subset of C stable under complex conjugation. Consider the algebra \(\mathcal{O}(S)\) of germs of complex-valued holomorphic functions in a neighborhood of S. The map \(f \mapsto f^{*}\) is an involution on \(\mathcal{O}(S)\) .

The subalgebra \(\mathcal{O}_{\mathbf{R}}(\mathrm{U})\) (resp. the subalgebra \(\mathcal{O}_{\mathbf{R}}(\mathrm{S})\) ) defined in no. 13 of I, p.~85 is the set of Hermitian elements of \(\mathcal{O}(\mathrm{U})\) (resp. of \(\mathcal{O}(\mathrm{S})\) ).

